\section{Model class and coordinate normalization}
\label{sec:potential}
On a fixed quasi-static state, assume a differentiable scalar landscape with
\begin{equation}
 \flux=\nabla_{\current}\coenergy,\qquad
 \coenergy(i_d,i_q)=\coenergy(i_d,-i_q).
 \label{eq:potential-gradient}
\end{equation}
The second condition is an additional reflection assumption. For a single
amplitude-invariant three-phase system $W'$ is physical co-energy divided by
$3/2$, so the gradient has the usual dq flux units.
Height is unobservable from slopes: $W'+C$ produces the same flux. Fix its
value at a supported reference as a gauge. The Hessian then supplies local
differential inductance, but the fit does not automatically establish
positive curvature or a true storage law in terminal-current coordinates.

Flux samples may be inconsistent with this class. Choosing a landscape whose
slopes best explain them is a projection under a basis, metric and regularizer.
Path integration is an alternative for exact conservative fields; with noisy
inconsistent data an arbitrary chosen path hides disagreement. The gradient
fit retains all sample directions and their residuals instead.
